% !TeX root = ../papers/nc-p3.tex
\section{階数14の斉次表示における全境界テンソル}
\label{sec:rank-fourteen-boundary}
\begin{theorem}\label{thm:rank-fourteen-boundary}
命題\ref{prop:rank-fourteen-descent}の斉次表示を固定し、二本の楕円曲線上で
四階テンソルの射影極限を取る。通常域の外の48幾何パラメータのうち、
次の三式の根の上の32点は標準Koszul左右AS正則で反対階数14である。
\[
 t^2-2t-7,\qquad t^2+6t+1,\qquad t^4+2t^2+24t+37 .
\]
残る16点の極限テンソルは標準14641型を与えない。
また通常な標数零の窓が非正則なら、その全ての101進還元で倍率座標は
$\infty,1,-1,-3,19,80$ のいずれかである。
\end{theorem}
\begin{proof}
薄い成分は元の表示のままにし、厚い二成分を
$\left(\begin{smallmatrix}A&X\\Y_0&Y_0A^{-1}X\end{smallmatrix}\right)$ と
$\left(\begin{smallmatrix}B&\lambda X\\Y_0/\lambda&Y_0B^{-1}X\end{smallmatrix}\right)^t$
で作る。ここで $X=\left(\begin{smallmatrix}1&1/\lambda\\1/\lambda&z\end{smallmatrix}\right)$、
$Y_0=\left(\begin{smallmatrix}1&\lambda\\\lambda z&1\end{smallmatrix}\right)$ である。
\path{computations/verify_rank_fourteen_boundary.py} はこれを楕円関数体で簡約する。
全成分の分母の既約因子は $t-1,t+1,t+3$ だけである。特に $A$ または $B$ が
特異になっても、積全体が正則ならその値から関係を回復できる。

通常倍率の非消滅条件を $t$ で因数分解すると、三つの極と無限遠以外では
定理の三式と $t^2+2t+5$ が尽くす。定理の三式は101で順に
\[
 t^2+99t+94,\quad t^2+6t+1,\quad
 (t^2+81t+40)(t^2+20t+59)
\]
となり、各二次式は既約である。各根で、必要なら剰余体を二次拡大し、
$i,Y$ の全四通りについて生のテンソルの隣接四切断と反対切断の階数を求める。
結果は $6,6,6,6$ と14で、各一文字切断は異なる文字に異なる列支持を持ち、階数4である。
四つの隣接切断から直接得た関係で、経路次元35,56、中心次元 $(3,2,0,0)$、
四中心商の長さ13までの列を全始点で確認した。

標数零でもテンソルの切断階数には一般階数による上界とこれらの剰余階数による
下界がある。従って関係空間は近傍でベクトル束をなし、一般正則点から全次数の経路次元下界を得る。
長さ4,5の余核が定階数となり、中心の下界と剰余体での上界も一致するので、
前定理の中心持上げと左右Koszul完全性の証明が適用される。
同じ議論はこの還元を持つ全ての通常な標数零の窓にも使える。
$t^2+2t+5$ の101での根は19と80であるから、前定理との合併で最後の主張を得る。

$t=-1+2i$ とその共役では $s=\pm1$ を代入し、$\mathbb Q(i)$ 上で計算すると
隣接四切断と反対切断は全て階数4である。最後の四基点は二次被覆の分岐点である。
有限の $a$ では $t-a=u^2$、無限遠では $t=u^{-2}$ と置く。
有理部分と $Y$ 倍の部分は位数の偶奇が異なる。全係数の最小位数を抜くと、
同じ位数の $Y$ の先頭係数は共通スカラーなので射影化で除ける。
次表を $\mathbb Q(i)$ 上の掃き出しで得る。
\[
\begin{array}{c|c|c}
 t & \text{四つの隣接階数} & \text{反対階数}\\\hline
 -3 & 2,1,2,1&2\\
 -1 & 2,2,2,2&1\\
  1 & 6,6,6,6&6\\
 \infty&2,2,2,2&1
\end{array}
\]
$t=1$ では三文字の関係挿入の余核次元が始点順に $28,28,22,22$ となり、
全て標準値20を超える。他の三基点と $t^2+2t+5$ の根では既に二次関係数が不足する。
非分岐の十個の $t$ では $i,Y$ の四組があり、分岐四点では各 $i$ に一つの極限となる。
これは指定した斉次表示を射影テンソル空間に延長したときの全境界を数えている。
\end{proof}
